\section{Finite generation of the multiplier}\label{sec:homology}

This section uses the arithmetic, strict-area comparison, and alternating
extension results of Lemmas~\ref{lem:arithmetic}, \ref{lem:comparison},
and~\ref{lem:extension}.  Its goal is
the following finiteness statement.  All homology groups in this Section have
integer coefficients unless another coefficient group is displayed.

\begin{theorem}\label{thm:multiplier}
There is an integer $m_{\mathrm{hom}}$ such that
$H_2(A_m;\ZZ)$ is finitely generated for every finite $m\geq m_{\mathrm{hom}}$.
\end{theorem}

The proof first establishes stability and finite generation for
$H_1(F_m)$ and $H_2(F_m)$.  Polygonal step functions enter through an explicit
resolution of a symmetric monoidal groupoid.  We then return to a fixed finite
number of tracks, prove that $F_m$ acts trivially on $H_2(A_m)$, and use the
homology spectral sequence of $A_m\triangleleft F_m$.  No presentation or
finite-generation assertion about $A_m$ is used in this Section.

This stability and group-completion strategy has precedents in the
Higman--Thompson calculation of Szymik and Wahl \cite{SzymikWahl2019}
and in Li's framework for ample groupoids and full groups
\cite{Li2025}.  Here the polygon category, its low-degree resolution,
and the return to $A_m$ on finitely many tracks are constructed below.
The external bar and group-completion results are cited at their
individual applications.

\subsection{Configurations of embedded squares}

An embedding $e:X\hookrightarrow mX$ means an injective partial piecewise
translation, defined on the entire Boolean square $X$.  Let $\mathcal E_m$ be
the semisimplicial set whose $p$-simplices are ordered tuples
$(e_0,\ldots,e_p)$ of embeddings with pairwise disjoint images.  Faces delete
entries.  Write $C_p(\mathcal E_m)$ for its integral chain group and augment
it by $C_{-1}(\mathcal E_m)=\ZZ$, sending every vertex to $1$.

\begin{lemma}\label{lem:hom-configuration-acyclic}
For each fixed integer $d\geq0$, the augmented complex
$C_*(\mathcal E_m)$ has zero homology in degrees $-1,0,\ldots,d$ when $m$ is
sufficiently large.  The required lower bound on $m$ depends only on $d$.
\end{lemma}

\begin{proof}
Any $b$ vertices have a common disjoint vertex if $m>b+1$: the union of their
images has area at most $b$, so its complement has area greater than $1$;
Lemma~\ref{lem:comparison} embeds $X$ in that complement.

We construct a chain contraction in the asserted degrees.  Choose a vertex
$v$ and put $h_{-1}(1)=v$.  Suppose $h_{-1},\ldots,h_{p-1}$ have been
constructed, satisfy $\partial h+h\partial=\id$ in degrees below $p$, and
the vertices appearing in $h_{p-1}(\rho)$ are bounded in number independently
of the $(p-1)$-simplex $\rho$.  For a $p$-simplex $\sigma$, the chain
\[
 z_\sigma=\sigma-h_{p-1}(\partial\sigma)
\]
is a cycle.  Its vertices have a bound depending only on $p$: if the bound
for $h_{p-1}$ is $b_{p-1}$, a valid bound is
$(p+1)+(p+1)b_{p-1}$.  Choose a vertex $w_\sigma$ disjoint from all these
vertices and define $h_p(\sigma)=w_\sigma*z_\sigma$, where the cone prepends
$w_\sigma$ to each ordered simplex.  The identity
\[
 \partial(w_\sigma*z)=z-w_\sigma*(\partial z)
\]
gives $\partial h_p(\sigma)+h_{p-1}(\partial\sigma)=\sigma$.
Moreover, $b_p=1+(p+1)+(p+1)b_{p-1}$ is a uniform bound for the vertices
appearing in $h_p(\sigma)$.  Starting with $b_{-1}=1$, choose $m$ larger
than all the finitely many bounds needed for $p\leq d$.  Linear extension
defines the required contraction.  There is no requirement that the
contraction be equivariant or that one vertex avoid an arbitrarily long
cycle.
\end{proof}

The group $F_m$ acts on $\mathcal E_m$ by postcomposition.  For a fixed $p$
and sufficiently large $m$, this action is transitive, even under $A_m$.
Indeed, the bijection from the first $p+1$ standard tracks to any given
configuration has area $p+1$ and extends by Lemma~\ref{lem:extension} when
$m>10(p+1)$.  The stabilizer in $F_m$ of the standard configuration fixes
those tracks pointwise and is therefore exactly $F_{m-p-1}$.  This avoids
any assertion that arbitrary equal-area polygons are equidecomposable.

\begin{lemma}\label{lem:hom-full-stability}
For $q=0,1,2$, adding an identity track induces an isomorphism
\[
 H_q(F_{m-1})\longrightarrow H_q(F_m)
\]
for all sufficiently large $m$.
The same assertion holds with $F_m$ replaced by the finite symmetric
group $\Sym(m)$.
\end{lemma}

\begin{proof}
Tensor the augmented configuration complex with a free right
$\ZZ[F_m]$-resolution of $\ZZ$ and take its total complex.  One filtration,
together with Lemma~\ref{lem:hom-configuration-acyclic}, shows that its
homology vanishes in any prescribed bounded range when $m$ is large.
The other filtration has, in columns $p\geq-1$,
\begin{equation}\label{eq:hom-configuration-ss}
 E^1_{p,q}=H_q(F_{m-p-1}),\qquad
 E^1_{-1,q}=H_q(F_m).
\end{equation}
Here and below only boundedly many columns are needed.  The identification
follows directly from the permutation module on simplices: restricting the
free resolution to a stabilizer is still a free resolution.

Every face inclusion of a standard stabilizer is a stabilization followed
by conjugation in the target stabilizer.  Inner conjugation induces the
identity on homology.  Thus $d^1$ is stabilization when $p$ is even and is
zero when $p$ is odd.  In particular, in row $q$ the entries at columns
$-1$ and $0$ of $E^2$ are respectively the cokernel and the kernel of
$H_q(F_{m-1})\to H_q(F_m)$.

The row $q=0$ is the alternating augmented row of copies of $\ZZ$ and is
exact.  For $q=1$, the only possible higher differentials into columns
$-1$ and $0$ come from row zero, so both entries vanish, as the total
homology vanishes.  This proves stability in degree one.  For $q=2$, the
possible sources for the entry $(-1,2)$ are $(1,1)$ and $(2,0)$, and those
for $(0,2)$ are $(2,1)$ and $(3,0)$.  The degree-one stability just proved
kills the indicated row-one entries when $m$ is increased by a fixed
amount; row zero is exact.  No higher differential has a nonnegative
source row, and there are no outgoing higher differentials from columns
$-1$ or $0$.  Hence the kernel and cokernel in degree two are zero.
Transitivity through column four and augmented acyclicity through degree
two suffice for this argument; all the requirements are finite.

For $\Sym(m)$ use ordered configurations of distinct points.  The same
bounded contraction works by choosing an unused point, the action is
transitive, and its stabilizer is $\Sym(m-p-1)$.  The identical spectral
sequence argument applies.
\end{proof}

\subsection{Group completion and three bars}

We now seek finite generation of the two stable homology groups in
Lemma~\ref{lem:hom-full-stability}.  Define $\mathcal B$ to be the
groupoid whose objects are finite lists $(U_1,\ldots,U_k)$ of Boolean
polygonal subsets of $X$.  A morphism is a piecewise translation bijection
between their disjoint unions, with the list entries serving as separate
tracks.  Tensor product is concatenation of lists; its symmetry reorders
the tracks.  The empty list is the unit.  Thus $\mathcal B$ is a strict
symmetric monoidal groupoid.  Let $\mathcal S$ be the groupoid of finite
sets and bijections, using its ordered skeleton with objects
$\{1,\ldots,n\}$, $n\geq0$.  Ordered disjoint union, with the second
block relabeled, makes its tensor product strictly associative and unital.

The groupoid $\mathcal S$ records the finite sets of track labels
over points of $X$.  By following these labels through strings of
exchanges, we will resolve $\mathcal B$ into levels built from finite
products of $\mathcal S$.  To compute their homology, we use three
bar constructions.  Reduced homology then starts in degree three,
so a product of two positive-degree classes first appears in degree
six.  The calculation through degree five is consequently additive
in the fibers; it is this range that will give stable homology
through degree two.

Here is a concrete bar convention, needed later for the resolution.
For a symmetric monoidal groupoid $\mathcal C$ and $p\geq0$, let
$D_p\mathcal C$ have objects
\[
 (C_{ij},\alpha_{ijk})_{0\leq i\leq j\leq k\leq p},\qquad
 \alpha_{ijk}:C_{ij}\sqcup C_{jk}\xrightarrow{\ \cong\ } C_{ik},
\]
with $C_{ii}$ the unit, the evident unit maps, and the associativity
condition at every four ordered indices.  Morphisms are families of
isomorphisms commuting with every $\alpha_{ijk}$.  Pullback along monotone
maps of index sets defines faces and degeneracies.  Tensor product is
coordinatewise, using the symmetry to interchange the middle summands in
the decomposition maps.  Evaluation on the elementary intervals gives an
equivalence
\begin{equation}\label{eq:hom-bar-product}
 D_p\mathcal C\simeq\mathcal C^p.
\end{equation}
An inverse replaces every interval by the ordered sum of its elementary
entries; coherence of the $\alpha$'s supplies its comparison isomorphism.

Write $B^0\mathcal C=|N\mathcal C|$ and let $B^d\mathcal C$ be the
realization of the nerve of
$D_{p_1}\cdots D_{p_d}\mathcal C$ over all $d$ bar indices.  We use
the usual realizations of these simplicial sets, or equivalently their
fat realizations.  Degeneracies are inclusions of simplicial subsets, so
these models are well based and levelwise equivalences remain
equivalences on realization.  For $M=|N\mathcal C|$, the first bar is
the usual $BM$: inserting ordered sums gives the level equivalences
with the ordinary simplicial bar.  The interval convention preserves
the symmetric product during iteration.

The coordinatewise tensor product and its symmetry make
$B^1\mathcal C$ a connected homotopy-commutative $H$-space.
We will use two standard bar facts, with their hypotheses visible in
this model.  The product equivalences
\eqref{eq:hom-bar-product}, the contractible level zero, and the
cofibration condition just noted satisfy the realization hypotheses in
\cite[Proposition~1.5 and its preceding Note; Proposition~A.1 and
Definition~A.4]{Segal1974}.
They imply that
\begin{equation}\label{eq:hom-deloop}
 B^{d-1}\mathcal C\simeq\Omega B^d\mathcal C\qquad(d\geq2).
\end{equation}
At these stages the preceding space is connected, so its component
monoid is already a group, as required by the delooping criterion.
The skeletal bar filtration also shows that
\begin{equation}\label{eq:hom-connectivity}
 \widetilde H_i(B^d\mathcal C)=0\quad(i<d),
 \qquad \pi_1(B^d\mathcal C)=0\quad(d\geq2).
\end{equation}
For example, after the first bar the spaces are connected.  In the next
bar the normalized level-zero row contributes only the base point, and
the first possible positive homology in the other levels gains one in
total degree.  The same product model, or the fundamental-group
calculation of the bar, gives simple connectivity from the second bar
onward.  This proves the displayed homological range inductively.

For $\mathcal C=\mathcal B$, our immediate target is therefore
finite generation of $H_i(B^3\mathcal B)$ for $i\leq5$.
Delooping will lower this bound to degree four for $B^2\mathcal B$
and then to degree three for $B^1\mathcal B$.  The spaces
$B^3\mathcal B$ and $B^2\mathcal B$ are simply connected, whereas
$B^1\mathcal B$ may have a nontrivial
fundamental group.  Its connected $H$-space structure will allow us
to pass to its universal cover without losing finite generation;
looping that cover gives the component $\Omega_0B^1\mathcal B$
of the constant loop, now through degree two.  Group completion
will identify this last homology with the stable homology of $F_m$.
The following lemma supplies the finiteness implications in this
chain.

\begin{lemma}\label{lem:hom-finiteness-transfers}
The following statements hold for spaces of CW homotopy type.
\begin{enumerate}[label=(\roman*)]
\item If $Z$ is simply connected, then, for every $r\geq0$,
$H_i(Z)$ is finitely generated for $i\leq r+1$ if and only if
$H_i(\Omega Z)$ is finitely generated for $i\leq r$.
\item Let $Y$ be a connected $H$-space.  If $H_i(Y)$ is finitely
generated for $i\leq r$, then $\pi_1(Y)$ and
$H_i(\widetilde Y)$ for $i\leq r$ are finitely generated, provided
$r\geq1$.  Conversely, if $\pi_1(Y)$ and those cover homology groups
are finitely generated, then $H_i(Y)$ is finitely generated for
$i\leq r$.
\end{enumerate}
\end{lemma}

\begin{proof}
We use the Serre homology spectral sequence in the form of
\cite[Theorem~1.3, p.~8]{HatcherSpectralSequences}.
For (i), apply it to the path-loop fibration, obtaining
\[
 E^2_{a,b}=H_a(Z;H_b(\Omega Z))\Longrightarrow H_{a+b}(PZ),
 \qquad PZ\simeq *.
\]
All coefficients are constant because $Z$ is simply connected.
If the base homology is finitely generated through degree $r+1$,
induct on $b\leq r$.  The term $E^2_{0,b}=H_b(\Omega Z)$ has no
outgoing differential.  Its incoming sources have strictly smaller
fiber degree and base degree at most $b+1$.  They are finitely
generated by induction and the universal coefficient theorem.  The
term at infinity is zero for $b>0$, so finitely many finitely
generated images exhaust $H_b(\Omega Z)$.  For the converse, induct
on $a\leq r+1$ in the row-zero term $E^2_{a,0}=H_a(Z)$.  It has no
incoming differential, its outgoing targets have smaller base
degree and fiber degree at most $a-1$, and its limiting term is
zero for $a>0$.  The successive quotients in this finite filtration
are subgroups of finitely generated groups.  This proves (i).

For (ii), the two products on based loops, concatenation and the
$H$-space product, give the same operation on $\pi_1(Y)$ and make
it abelian.  Thus $\pi_1(Y)=H_1(Y)$.  This fundamental group acts
trivially on the homology of the universal cover.  To see the latter
claim, represent a deck transformation by a loop $\alpha$ at the
unit.  Lift the family of maps
$(t,y)\mapsto\alpha(t)y$ to the universal cover, starting with the
identity (using the unit homotopy if necessary).  Its terminal map
is the deck transformation represented by $\alpha$.  Consequently
that map is homotopic to the identity.

Put $Q=\pi_1(Y)$.  The spectral sequence for the universal cover is
\[
 E^2_{a,b}=H_a(Q;H_b(\widetilde Y))\Longrightarrow H_{a+b}(Y),
\]
with constant coefficients.  A finitely generated abelian group
has finitely generated homology in each degree with finitely
generated constant coefficients: tensor the resolutions for an
infinite cyclic group and for a finite cyclic group.  Induction on
$b$ in column zero now proves the forward implication, since all
incoming sources have lower fiber degree and the limiting term is
a subgroup of $H_b(Y)$.  The reverse implication follows directly
from the finitely many terms on each total-degree diagonal.
\end{proof}

\begin{lemma}\label{lem:hom-cofinal-completion}
There are natural homology isomorphisms
\begin{align*}
 \underset{m}{\operatorname{colim}}\,H_j(F_m)
   &\cong H_j(\Omega_0 B^1\mathcal B),\\
 \underset{m}{\operatorname{colim}}\,H_j(\Sym(m))
   &\cong H_j(\Omega_0 B^1\mathcal S).
\end{align*}
Here $\Omega_0$ denotes the component of the constant loop.
\end{lemma}

\begin{proof}
The monoid $M=|N\mathcal B|$ is well based and has CW homotopy type.
The symmetry supplies a homotopy between the two orders of
multiplication, so $\pi_0(M)$ is central in the integral Pontryagin
ring.  The integral group-completion theorem
\cite[Proposition~1]{McDuffSegal1976} therefore identifies the
homology of $\Omega BM$ with the localization of $H_*(M)$ at its
components.

There is one cofinal object here.  The object
$U=(U_1,\ldots,U_k)$ has complement
$U^c=(X\setminus U_1,\ldots,X\setminus U_k)$ and
$U\sqcup U^c\cong kX$.  Hence inverting $X$ inverts every
component.  Its localization is computed by the telescope
\[
 T=\operatorname{tel}(M\xrightarrow{X\sqcup-}M
                  \xrightarrow{X\sqcup-}\cdots).
\]
Equivalently, every right translation on this telescope is a
homology equivalence: translation by $X$ is invertible, tensor
symmetry exchanges left and right translations, and a complement
gives an inverse for translation by any object.  This verifies the
action hypothesis in the telescope form of group completion \cite[Proposition~2 and p.~281]{McDuffSegal1976};
see also the corrected proof in
\cite[Theorem~4.1]{MillerPalmer2015}, applied with ordinary integral
homology.  In the nerve-realization model, $M$ is a Hausdorff,
locally contractible CW complex and its telescope is Hausdorff,
as required there.

The bar two-skeleton identifies $\pi_1(BM)$ with the group
completion of the commutative monoid of object-isomorphism classes:
an object supplies a loop, and a pair supplies the relation that
its sum represents the product.  The component of an object $U$
at stage $r$ of $T$ is therefore $[U]-r[X]$.  If this is zero,
there is an object $V$ with
$U\sqcup V\cong rX\sqcup V$.  Choose $V'$ with
$V\sqcup V'\cong kX$.  Then
\[
 U\sqcup kX\cong(r+k)X.
\]
Thus every component occurring in the neutral telescope component
eventually becomes a standard component $BF_{r+k}$.  Choices of
the identifying isomorphism differ by an inner automorphism and
give the same map on homology.  It follows that the neutral
component has homology $\operatorname{colim}_m H_j(F_m)$, with
the usual stabilization maps.  Group completion respects these
component gradings, proving the first assertion.  For
$\mathcal S$, use a singleton in place of $X$; the same argument
gives the second assertion.
\end{proof}

\begin{lemma}\label{lem:hom-finite-set-bars}
For $d=1,2,3$, the groups $H_i(B^d\mathcal S)$ are finitely
generated for $i\leq d+2$.
\end{lemma}

\begin{proof}
Finite groups have finitely generated integral homology in every
degree, as their bar chain groups have finite rank in each degree.
Lemma~\ref{lem:hom-full-stability} therefore gives finite
generation of the stable homology of $\Sym(m)$ through degree two.
By Lemma~\ref{lem:hom-cofinal-completion}, this is the homology
through degree two of $\Omega_0 B^1\mathcal S$, or equivalently
of $\Omega\widetilde{B^1\mathcal S}$.
Lemma~\ref{lem:hom-finiteness-transfers}(i) gives finite
generation for the cover through degree three.  The fundamental
group of $B^1\mathcal S$ is the group completion of $\NN$, hence
$\ZZ$.  Part (ii) gives finite generation for $B^1\mathcal S$
through degree three.

Ascend the bars using their skeletal filtrations and the product
models \eqref{eq:hom-bar-product}.  In the normalized filtration
for the next bar, column zero is a point, so a positive-degree
term of total degree at most $d+2$ has bar degree at least one
and internal degree at most $d+1$.  The induction hypothesis and
the integral K\"unneth theorem make the homology of every finite
product in this range finitely generated.  There are finitely
many bidegrees on each relevant diagonal.  This proves the
assertion successively for $d=2$ and $d=3$.
\end{proof}

Consequently the three coefficient groups
\begin{equation}\label{eq:hom-Kj}
 K_j=H_j(B^3\mathcal S),\qquad j=3,4,5,
\end{equation}
are finitely generated.  The remaining task is to obtain the same
degree-five bound for $B^3\mathcal B$.  We first describe its
resolution completely, including the maps that identify it.

\subsection{Strings and the trajectory differential}

We now construct the resolution whose levels have finite-set fibers.
A string of exchanges records both the positions visited by a point
and the track labels carried along that path.  Keeping these two
kinds of data separate will turn its low-degree homology into a
translation-homology calculation.

A position-preserving isomorphism of objects of $\mathcal B$ may
change the track label, piecewise, but leaves the point of $X$
unchanged.  Denote the subgroupoid of such isomorphisms, with all
objects retained, by $\mathcal P$.  For $p\geq0$ define
$\mathcal C_p$ as follows.  Its objects are strings
\[
 U_0\xrightarrow{f_1}U_1\xrightarrow{f_2}\cdots
       \xrightarrow{f_p}U_p
\]
of arbitrary morphisms of $\mathcal B$.  Its morphisms are
commuting ladders whose maps $U_i\to U'_i$ belong to
$\mathcal P$.  Deleting a vertex and composing adjacent arrows
defines the faces; inserting an identity defines the degeneracies.
Concatenation gives a simplicial symmetric monoidal groupoid.
A simplex of $N_q\mathcal C_p$ is a commuting grid with $p$
horizontal arrows and $q$ vertical arrows: the horizontal arrows
are arbitrary exchanges, and the vertical arrows preserve positions.

We first describe these strings by their trajectories and finite fibers.
For $x\in X$ and $\gamma_1,\ldots,\gamma_p\in\Gamma$, a
trajectory is the sequence of positions
\[
 x,\quad x+\gamma_1,\quad x+\gamma_1+\gamma_2,\quad\ldots,
 \quad x+\gamma_1+\cdots+\gamma_p.
\]
Cut a string in $\mathcal C_p$ according to its successive
translation charts and the inverse images of later cuts.  Each
resulting strand has a fixed increment tuple
$(\gamma_1,\ldots,\gamma_p)$ and a Boolean set of initial points.
The translation action on $X$ is free, so the increments are
uniquely determined by the positions.  Track multiplicities give
a finite set over each initial point and increment tuple.
Position-preserving ladders are exactly permutations of these
fibers.  Conversely, a Boolean locally constant family of finite
sets on $X\times\Gamma^p$, supported on finitely many increment
tuples, reconstructs a string by taking its strands as separate
list entries and translating them successively.  These
constructions give an equivalence of symmetric monoidal
groupoids.

More explicitly, restrict to a finite set of increment tuples
and a finite Boolean partition on which all fiber sets and maps
are constant.  The corresponding groupoid is a finite product
of copies of $\mathcal S$.  Enlarging the finite support inserts
empty fibers; refining an atom duplicates its fiber.  Every
finite diagram of strings and isomorphisms occurs at one such
stage.  Thus nerves, bars, and their chain complexes are obtained
by a filtered colimit of these finite stages.

The levels are now expressed in terms of finite sets.  The next lemma
shows that realizing the spaces $B^3\mathcal C_p$ in the string index
$p$ recovers $B^3\mathcal B$.  The commuting-grid description is the
basis of its proof.

\begin{lemma}\label{lem:hom-string-resolution}
There is a zigzag of natural realization equivalences between
$B^3\mathcal B$ and the realization in $p$ of $B^3\mathcal C_p$.
\end{lemma}

\begin{proof}
Fix three bar indices $\boldsymbol n=(n_1,n_2,n_3)$, and put
$\mathcal D=D_{n_1}D_{n_2}D_{n_3}\mathcal B$.
Let $\mathcal D^{\mathrm{pos}}$ have those objects of
$\mathcal D$ whose decomposition isomorphisms at every bar level
preserve positions in $X$, but retain \emph{all} horizontal diagram isomorphisms
between them.  The inclusion
\begin{equation}\label{eq:hom-pos-inclusion}
 \mathcal D^{\mathrm{pos}}\longrightarrow\mathcal D
\end{equation}
is full and essentially surjective.  Indeed, replace every
multi-interval object by the lexicographically ordered sum of its
elementary entries, indexed by products of consecutive intervals
$[i_1-1,i_1]\times[i_2-1,i_2]\times[i_3-1,i_3]$ in the three
bar-index sets.  Its decomposition maps are
canonical reorderings and preserve positions.  The original
coherent decomposition maps identify each ordered sum with the
corresponding original multi-interval.  They form a diagram
isomorphism because changing the order of the cuts only introduces
the specified symmetric reordering.  The comparison is allowed to
translate pieces: it is a horizontal arrow.  For example, in one
bar at degree two, the long edge $U_{02}$ is replaced by
$U_{01}\sqcup U_{12}$, and the comparison on that edge is the
original map $\alpha_{012}$, regardless of whether it preserves
positions.

The inclusions \eqref{eq:hom-pos-inclusion} commute with all bar
operations.  A face pulls back the existing interval diagram and its
specified decomposition maps; it does not choose a new ordered-sum
normalization.  Composing position-preserving decompositions,
inserting a unit, and reordering sums remain position-preserving.
We use these inclusions as the natural maps; no strictly natural
choice of an inverse normalization is required.

Now take the double nerve whose horizontal strings use arbitrary
arrows of $\mathcal D^{\mathrm{pos}}$ and whose vertical arrows
preserve positions.  Reading it horizontally first gives exactly
the nerve of $D_{n_1}D_{n_2}D_{n_3}\mathcal C_p$: its
decomposition maps are position-preserving ladders, precisely the
condition built into the morphisms of $\mathcal C_p$.

Read the same double nerve vertically first, fixing a vertical
string length $q$.  Its groupoid has objects strings
$V_0\xrightarrow{v_1}\cdots\xrightarrow{v_q}V_q$ of
position-preserving arrows and has arbitrary horizontal commuting
ladders.  Evaluation at $V_0$ is an equivalence with
$\mathcal D^{\mathrm{pos}}$.  It is essentially surjective by
constant strings.  It is fully faithful because a horizontal
arrow $a_0:V_0\to W_0$ uniquely determines all its components:
\begin{equation}\label{eq:hom-ladder-formula}
 a_i=w_{0i}\,a_0\,v_{0i}^{-1},
 \qquad v_{0i}=v_i\cdots v_1,
 \quad w_{0i}=w_i\cdots w_1.
\end{equation}
Every map in \eqref{eq:hom-ladder-formula} is an allowed
horizontal diagram map.  Constant-string inclusions are natural
in $q$ and in every bar index, and are levelwise equivalences.
Thus the constant-string maps and
\eqref{eq:hom-pos-inclusion} give the asserted zigzag after taking
nerves and successive realizations.  The levelwise-equivalence
criterion applies to these simplicial-set models as explained
before \eqref{eq:hom-deloop}.
\end{proof}

We now apply this finite-fiber description to homology.
Put $E=B^3\mathcal S$.  By \eqref{eq:hom-connectivity},
$H_0(E)=\ZZ$ and $\widetilde H_i(E)=0$ for $i<3$.
For a finite product $E^s$, the integral K\"unneth formula gives
\begin{equation}\label{eq:hom-product-additivity}
 H_j(E^s)=\bigoplus_{a=1}^s H_j(E)
 \qquad(j=3,4,5).
\end{equation}
A tensor term using two positive-degree factors first has degree
six; a corresponding $\operatorname{Tor}$ term first has degree
seven.  Terms involving $H_0(E)=\ZZ$ have no such torsion term.
The disjoint-union map on the fibers restricts to the identity
on either factor when the other is empty.  Hence under
\eqref{eq:hom-product-additivity} it induces ordinary addition.
Refinement induces the diagonal, since the same fiber is copied
onto every new atom.  Exactness of filtered colimits therefore
identifies
\begin{equation}\label{eq:hom-trajectory-chains}
 H_j(B^3\mathcal C_p)
  =\bigoplus_{(\gamma_1,\ldots,\gamma_p)\in\Gamma^p}
        C(X,\ZZ)\otimes K_j\qquad(j=3,4,5),
\end{equation}
where $C(X,\ZZ)$ denotes the Boolean step functions.  There is
no derived tensor hidden here: $C(X,\ZZ)$ is the filtered union
of the free groups of functions on finite Boolean partitions.

For clarity, all the differential maps in
\eqref{eq:hom-trajectory-chains} can be written on trajectories:
\begin{align}
 d_0(x;\gamma_1,\ldots,\gamma_p)
   &=(x+\gamma_1;\gamma_2,\ldots,\gamma_p),\notag\\
 d_i(x;\gamma_1,\ldots,\gamma_p)
   &=(x;\gamma_1,\ldots,\gamma_i+\gamma_{i+1},\ldots,\gamma_p),
       &&0<i<p,\label{eq:hom-trajectory-faces}\\
 d_p(x;\gamma_1,\ldots,\gamma_p)
   &=(x;\gamma_1,\ldots,\gamma_{p-1}).\notag
\end{align}
Fibers acquiring the same label are disjointly summed.  In the
coefficient convention $\gamma_*f(y)=f(y-\gamma)$, the first face
therefore applies $\gamma_{1*}$ to the coefficient, each interior
face adds two consecutive increments, and the last face leaves
the coefficient unchanged.  The alternating sum of these maps is
the inhomogeneous group-homology bar differential for the
translation module $C(X,\ZZ)\otimes K_j$ (the group $\Gamma$ is
abelian, so the same pushforward convention also gives a right
module).

The skeletal spectral sequence for the resolution in
Lemma~\ref{lem:hom-string-resolution} consequently has, in the
reduced rows relevant to total degree at most five,
\begin{equation}\label{eq:hom-trajectory-ss}
 E^2_{p,j}
   =H_p\bigl(\Gamma;C(X,\ZZ)\otimes K_j\bigr),
       \qquad j=3,4,5.
\end{equation}
Rows one and two are zero.  Row zero is the constant simplicial
group $\ZZ$, since every $B^3\mathcal C_p$ is connected; its
homology contributes only the degree-zero copy of $\ZZ$.
We now establish finite generation of the coefficient homology
in \eqref{eq:hom-trajectory-ss}.

\subsection{Corner symbols and polygonal coefficients}

Let $H=\OO^2$, regarded as a discrete free abelian group of rank
four, and let $L$ be the abelian group of compactly supported
integer-valued step functions on the plane whose boundaries use
finitely many allowable lines.  Equalities are modulo sets with
empty interior.  Translations make $L$ a module over
\[
 R=\ZZ[H]\cong
       \ZZ[t_1^{\pm1},t_2^{\pm1},t_3^{\pm1},t_4^{\pm1}],
\]
a noetherian ring.

Translation modules of polyhedral step functions also occur in the
homology of projection-method patterns
\cite[Section~3.1]{GahlerHuntonKellendonk2013}.  Our immediate task is
the integral finite-generation statement below.  We prove it by an
explicit corner map and the finite vertex types of
Lemma~\ref{lem:arithmetic}, without a cut-and-project identification.

\begin{lemma}\label{lem:hom-corner-module}
The $R$-module $L$ is finitely generated.
\end{lemma}

\begin{proof}
Let $V$ be the set of intersections of at least two nonparallel
allowable lines.  At $z\in V$, let $r(z)\in\{2,3,4\}$ be the
number of allowable lines through $z$.  Their $2r(z)$ sectors
give a free abelian group of sector germs.  Quotient it by the
constant germ and the indicator germs of either half-plane for
each incident line; call the quotient $Q_z$.

We spell out this integral quotient.  Number the sectors
cyclically, put $r=r(z)$, and write a germ as
$(a_0,\ldots,a_{2r-1})$.  With indices modulo $2r$, let
\[
 d_i=a_{i+1}-a_i,\qquad q_i=d_i+d_{i+r}\quad(0\leq i<r).
\]
Then $\sum_iq_i=0$, and
\begin{equation}\label{eq:hom-corner-quotient}
 Q_z\cong\{(q_0,\ldots,q_{r-1})\in\ZZ^r:\textstyle\sum_iq_i=0\}
       \cong\ZZ^{r-1}.
\end{equation}
Indeed, cyclic differences identify germs modulo constants with
$\{d\in\ZZ^{2r}:\sum_i d_i=0\}$.  A half-plane germ has
difference vector $\pm(e_i-e_{i+r})$.  These vectors generate
exactly the kernel of the pair-sum map $d\mapsto q$; the map is
onto, since $d=(q,0)$ has zero total sum.  This proves
\eqref{eq:hom-corner-quotient} over $\ZZ$.

For $f\in L$, record its sector germ modulo this subgroup at
every $z$.  The result is a translation-equivariant map
\[
 c:L\longrightarrow\bigoplus_{z\in V}Q_z.
\]
This has finite support.  To define the germ, use the finite line
family defining $f$ and take a small neighborhood excluding its
nonincident lines; other allowable directions through $z$
merely refine sectors with equal values.  A nonzero symbol can
occur only where two of the finitely many defining lines cross.
The density of the collection of all allowable lines causes no
difficulty.

The map $c$ is injective.  If $c(f)=0$, the equation
$d_{i+r}=-d_i$ says that the jump across each incident oriented
line is the same on both sides of the vertex.  Fix a line in a
finite defining family for $f$.  Its jump is constant between
successive intersections with that family and, by the preceding
identity, continues unchanged across every intersection.  It is
therefore constant along the whole line.  Compact support makes
the jump zero far away, hence everywhere.  All the line jumps
vanish.  Any two chambers of the finite line arrangement can be
joined by a path crossing its lines away from vertices; their
function values consequently agree.  The common value is zero
outside the compact support, so $f=0$.

Lemma~\ref{lem:arithmetic} places $V$ inside
$D^{-1}\OO^2$ for a fixed integer $D>0$.  Thus $V$ has finitely
many $H$-orbits; translation preserves the full set of incident
directions, and a point has trivial translation stabilizer.
After choosing one representative from each orbit,
\[
 \bigoplus_{z\in V}Q_z
       \cong\bigoplus_{\alpha=1}^s
                    R^{\,r(z_\alpha)-1}.
\]
This is a finite free $R$-module.  Its submodule $c(L)$ is
finitely generated because $R$ is noetherian, proving the Lemma.
\end{proof}

\begin{lemma}\label{lem:hom-coefficient-finiteness}
For every finitely generated abelian group $K$ and every $i\geq0$,
\[
 H_i\bigl(\Gamma;C(X,\ZZ)\otimes K\bigr)
\]
is a finitely generated abelian group.
\end{lemma}

\begin{proof}
Put $N=\ZZ^2$ and $J=\tau\ZZ^2$.  Then
$H=N\oplus J$ and $J\cong\Gamma$.  Cutting a compactly
supported function along the integer square grid gives, as
$N$-modules,
\begin{equation}\label{eq:hom-square-induction}
 L\cong\ZZ[N]\otimes C(X,\ZZ).
\end{equation}
The second factor is represented by functions supported in one
unit square, and $N$ shifts the square index.  Only finitely many
squares meet any given support.  The Boolean convention makes
the half-open choice on square edges immaterial.

Tensor \eqref{eq:hom-square-induction} with $K$.  Induction from
the trivial subgroup, or the two-generator free abelian
resolution, gives
\[
 H_q(N;L\otimes K)=
 \begin{cases}
 C(X,\ZZ)\otimes K,&q=0,\\
 0,&q>0.
 \end{cases}
\]
On the coinvariants, a translation in $J$ cuts and moves pieces
back to the chosen square, so its action is exactly translation
modulo one.  Taking the two successive free abelian resolutions
for $N$ and $J$ therefore identifies
\begin{equation}\label{eq:hom-plane-torus}
 H_i(H;L\otimes K)
   \cong H_i\bigl(\Gamma;C(X,\ZZ)\otimes K\bigr).
\end{equation}

By Lemma~\ref{lem:hom-corner-module}, $L\otimes K$ is a
finitely generated $R$-module: tensor a finite set of $R$-module
generators of $L$ with a finite set of abelian generators of $K$.
The Koszul resolution on $t_1-1,\ldots,t_4-1$ computes the
left side of \eqref{eq:hom-plane-torus}.  Its chain groups are
finite sums of $L\otimes K$, and its homology is finitely
generated over $R$ by noetherianity.  On this complex, exterior
multiplication $\varepsilon_a$ by the $a$th basis vector satisfies
\[
 d\varepsilon_a+\varepsilon_a d=(t_a-1)\id.
\]
Thus each $t_a-1$ acts trivially on homology.  The homology is
therefore a finitely generated module over
$R/(t_1-1,\ldots,t_4-1)=\ZZ$, as required.  We have proved
finite generation of $L$ integrally before tensoring; no
preservation of the corner-map injection under tensor product
is being assumed.
\end{proof}

\begin{proposition}\label{prop:hom-full-finiteness}
For all sufficiently large finite $m$, the groups $H_1(F_m)$
and $H_2(F_m)$ are finitely generated.
\end{proposition}

\begin{proof}
By Lemma~\ref{lem:hom-finite-set-bars}, each $K_j$ in
\eqref{eq:hom-Kj} is finitely generated.
Lemma~\ref{lem:hom-coefficient-finiteness} makes every relevant
term of \eqref{eq:hom-trajectory-ss} finitely generated.
There are finitely many terms of total degree at most five,
apart from the rows already shown to vanish.  Hence
Lemma~\ref{lem:hom-string-resolution} gives finite generation
of $H_i(B^3\mathcal B)$ for $i\leq5$.

Apply \eqref{eq:hom-deloop} and
Lemma~\ref{lem:hom-finiteness-transfers}(i) twice.  First
$H_i(B^2\mathcal B)$ is finitely generated for $i\leq4$, and
then $H_i(B^1\mathcal B)$ is finitely generated for $i\leq3$.
Both bases in these applications are simply connected by
\eqref{eq:hom-connectivity}.  The connected $H$-space
$B^1\mathcal B$ has finitely generated abelian fundamental
group, and part (ii) of the same Lemma gives finite generation
of its universal-cover homology through degree three.  Part (i)
now gives finite generation of
$H_i(\Omega_0 B^1\mathcal B)$ for $i\leq2$.
Lemma~\ref{lem:hom-cofinal-completion} identifies these with
the stable homology groups of $F_m$, and
Lemma~\ref{lem:hom-full-stability} transfers the conclusion to
every sufficiently large finite $m$.
\end{proof}

\subsection{Returning to the alternating group on finitely many tracks}

Only the passage from full-group homology to the multiplier
remains.  We must control the conjugation action on the whole
group $H_2(A_m)$, because a bound on its coinvariants alone would
not prove Theorem~\ref{thm:multiplier}.

\begin{lemma}\label{lem:hom-alternating-surjection}
For all sufficiently large $m$, stabilization induces a
surjection $H_2(A_{m-1})\to H_2(A_m)$.
\end{lemma}

\begin{proof}
Use the action of $A_m$ on $\mathcal E_m$.  It is transitive
in the bounded range of degrees needed here, by
Lemma~\ref{lem:extension}.  For the standard $p$-simplex, its
stabilizer is
\[
 A_m\cap F_{m-p-1}
 =\ker\bigl(F_{m-p-1}\longrightarrow F_m/A_m\bigr).
\]
By Theorem~\ref{thm:simplicity}, $A_t=[F_t,F_t]$.
The degree-one stability from
Lemma~\ref{lem:hom-full-stability} identifies
$H_1(F_{m-p-1})\to H_1(F_m)$ as an isomorphism when $m$ is
large and $p$ is bounded.  The displayed kernel is therefore
exactly $A_{m-p-1}$.

The augmented spectral sequence now has
$E^1_{p,q}=H_q(A_{m-p-1})$ and augmented column
$E^1_{-1,q}=H_q(A_m)$.  Its $E^2_{-1,2}$ is the cokernel
of the stabilization map in the statement.  Every $A_t$ is
perfect, so the row $q=1$ vanishes; row zero is the exact
alternating row of copies of $\ZZ$.  The only possible higher
incoming differentials at $(-1,2)$ have sources $(1,1)$ for
$d^2$ and $(2,0)$ for $d^3$.  They are zero on their respective
pages.  There is no outgoing differential from column $-1$.
The limiting term vanishes by
Lemma~\ref{lem:hom-configuration-acyclic}, so the cokernel
is zero.
\end{proof}

\begin{lemma}\label{lem:hom-trivial-action}
For all sufficiently large finite $m$, conjugation by every
element of $F_m$ induces the identity on $H_2(A_m)$.
\end{lemma}

\begin{proof}
Choose one finite integer $b$ such that all the stabilization
maps in Lemma~\ref{lem:hom-alternating-surjection} are
surjective once the target index exceeds $b$.  Composition gives
\[
 H_2(A_b)\twoheadrightarrow H_2(A_m)\qquad(m\geq b).
\]
Let $Y$ be the first $b$ tracks of $mX$, and take $m>10b$.
For $f\in F_m$, the restriction $f|_Y:Y\to f(Y)$ is a
piecewise translation between sets of area $b$.
Lemma~\ref{lem:extension} supplies $a\in A_m$ that agrees
with $f$ on $Y$.  For every $h\in A_b$, viewed as supported
on $Y$, one has the exact equality
\[
 f h f^{-1}=a h a^{-1}.
\]
Conjugation by $a$ is inner in $A_m$ and acts trivially on
its homology.  The two induced maps agree on the image of
$H_2(A_b)$, which is all of $H_2(A_m)$.  This proves the
claim.  The bank $Y$ is fixed independently of both $f$ and the
homology class; only the extending element $a$ depends on $f$.
\end{proof}

\begin{proof}[Proof of Theorem~\ref{thm:multiplier}]
Fix a finite $m$ beyond the thresholds in
Proposition~\ref{prop:hom-full-finiteness} and
Lemma~\ref{lem:hom-trivial-action}.  Put $Q=F_m/A_m$.
Theorem~\ref{thm:simplicity} identifies $Q$ with $H_1(F_m)$,
so $Q$ is finitely generated abelian.  Apply the integral
Lyndon--Hochschild--Serre spectral sequence
\cite[6.8.2, pp.~195--196]{Weibel1994} to
$1\to A_m\to F_m\to Q\to1$:
\[
 E^2_{p,q}=H_p(Q;H_q(A_m))\Longrightarrow H_{p+q}(F_m).
\]
By Lemma~\ref{lem:hom-trivial-action},
$E^2_{0,2}=H_2(A_m)$ itself.  Perfection gives
$H_1(A_m)=0$, so the possible incoming $d^2$ from $(2,1)$
vanishes.  The only remaining possible incoming differential is
\[
 d^3:E^3_{3,0}\longrightarrow E^3_{0,2}=H_2(A_m).
\]
Its source is a subquotient of $H_3(Q)$ and is finitely
generated.  Later incoming differentials have negative source
row, and none can leave column zero.  Consequently
\[
 E^\infty_{0,2}=H_2(A_m)/\im(d^3).
\]
This limiting term is the bottom filtration subgroup of
$H_2(F_m)$, which is finitely generated by
Proposition~\ref{prop:hom-full-finiteness}.  Both
$\im(d^3)$ and the displayed quotient are finitely generated,
so $H_2(A_m)$ is finitely generated.  All thresholds used in
choosing $m$ are fixed finite integers, and the argument applies
to every larger finite $m$.
\end{proof}
